% ============================================================
% DOCUMENT 1 : CONCEPT NOTE — eSchool Ecosystem
% Titre : "eSchool Ecosystem : Vision et Stratégie"
% Public : Décideurs, investisseurs, partenaires éducatifs
% ============================================================
\documentclass[12pt,a4paper,oneside]{report}

% --- Encodage et langue ---
\usepackage[utf8]{inputenc}
\usepackage[T1]{fontenc}
\usepackage[french]{babel}
\usepackage{lmodern}

% --- Mise en page ---
\usepackage[margin=2.5cm]{geometry}
\usepackage{setspace}
\onehalfspacing
\usepackage{parskip}
\usepackage{fancyhdr}

% --- Graphiques et couleurs ---
\usepackage{graphicx}
\usepackage[dvipsnames,svgnames,x11names]{xcolor}
\usepackage{tikz}
\usetikzlibrary{shapes.geometric, arrows.meta, positioning, calc, fit, backgrounds, shadows, decorations.pathmorphing}

% --- Tableaux ---
\usepackage{booktabs}
\usepackage{tabularx}
\usepackage{longtable}
\usepackage{multirow}
\usepackage{array}

% --- Code et listings ---
\usepackage{listings}
\lstset{
  basicstyle=\ttfamily\small,
  backgroundcolor=\color{gray!8},
  frame=single,
  rulecolor=\color{gray!30},
  breaklines=true,
  captionpos=b,
  tabsize=2
}

% --- Boîtes colorées ---
\usepackage[most]{tcolorbox}
\newtcolorbox{infobox}[1][]{colback=blue!5,colframe=blue!50!black,fonttitle=\bfseries,title=#1}
\newtcolorbox{warnbox}[1][]{colback=orange!5,colframe=orange!70!black,fonttitle=\bfseries,title=#1}
\newtcolorbox{successbox}[1][]{colback=green!5,colframe=green!50!black,fonttitle=\bfseries,title=#1}
\newtcolorbox{conceptbox}[1][]{colback=purple!5,colframe=purple!50!black,fonttitle=\bfseries,title=#1}

% --- Hyperliens ---
\usepackage[colorlinks=true,linkcolor=blue!60!black,urlcolor=blue!50!black,citecolor=green!50!black]{hyperref}
\usepackage{bookmark}

% --- En-têtes et pieds de page ---
\pagestyle{fancy}
\fancyhf{}
\fancyhead[L]{\small\leftmark}
\fancyhead[R]{\small eSchool Ecosystem}
\fancyfoot[C]{\thepage}
\renewcommand{\headrulewidth}{0.4pt}

% --- Couleurs personnalisées ---
\definecolor{eschoolblue}{HTML}{2563EB}
\definecolor{eschoolpurple}{HTML}{7C3AED}
\definecolor{eschoolteal}{HTML}{0D9488}
\definecolor{eschoolamber}{HTML}{D97706}
\definecolor{hubcolor}{HTML}{3B82F6}
\definecolor{spokecolor}{HTML}{10B981}
\definecolor{arrowcolor}{HTML}{6366F1}

% ============================================================
\begin{document}

% --- PAGE DE COUVERTURE ---
\begin{titlepage}
\begin{tikzpicture}[remember picture, overlay]
  % Fond gradient
  \fill[left color=eschoolblue, right color=eschoolpurple] 
    (current page.south west) rectangle (current page.north east);
  % Cercles décoratifs
  \foreach \x/\y/\r/\o in {3/24/4/0.08, 16/20/3/0.06, 5/6/5/0.05, 18/8/2.5/0.1} {
    \fill[white, opacity=\o] (\x,\y) circle (\r);
  }
  % Bande blanche en bas
  \fill[white, opacity=0.95] 
    ([yshift=6cm]current page.south west) rectangle (current page.south east);
\end{tikzpicture}

\vspace*{4cm}
\begin{center}
  {\Huge\bfseries\color{white} eSchool Ecosystem}\\[0.8cm]
  {\LARGE\color{white!90} Vision et Stratégie}\\[0.5cm]
  {\Large\color{white!80} Concept Note}\\[3cm]
  \rule{8cm}{0.5pt}\\[0.5cm]
  {\large\color{white!85} Plateforme Intégrée de Gestion}\\
  {\large\color{white!85} de l'Éducation en Afrique de l'Est}\\[4cm]
\end{center}
\begin{center}
  {\normalsize\color{gray!80} Version 1.0 — Mars 2026}\\[0.3cm]
  {\normalsize\color{gray!80} Projet eSchool / MonEcole — evanet08}\\[0.3cm]
  {\normalsize\color{gray!80} \texttt{https://eschool-rdc.pro}}
\end{center}
\end{titlepage}

% --- TABLE DES MATIÈRES ---
\tableofcontents
\newpage

% ============================================================
% CHAPITRE 1 : RÉSUMÉ EXÉCUTIF
% ============================================================
\chapter{Résumé Exécutif}

\section{Vision}

L'écosystème \textbf{eSchool} est une plateforme numérique intégrée conçue pour transformer la gestion éducative en Afrique de l'Est et Centrale. Face à la fragmentation des outils existants et à l'absence de standards numériques unifié, eSchool propose une architecture \textbf{Hub-and-Spoke} qui sépare intelligemment la gouvernance nationale de la gestion opérationnelle scolaire.

\begin{conceptbox}[Proposition de Valeur]
\textbf{Un ministère structure, mille écoles exécutent.} \\
eSchool offre une solution à \textbf{deux niveaux} : un hub centralisé (\texttt{eSchoolStructure}) pour les standards nationaux, et des instances décentralisées (\texttt{MonEcole}) pour chaque établissement scolaire, le tout synchronisé en temps réel via des vues MySQL cross-bases.
\end{conceptbox}

\section{Les Deux Plateformes}

\begin{center}
\begin{tikzpicture}[
  block/.style={rectangle, rounded corners=8pt, draw, text width=6cm, minimum height=3cm, align=center, font=\small, drop shadow},
  arrow/.style={-{Stealth[length=8pt]}, thick, arrowcolor}
]
  % Hub
  \node[block, fill=hubcolor!15, draw=hubcolor!60] (hub) {
    \textbf{\large eSchool Structure}\\[4pt]
    \textcolor{hubcolor}{\textbf{Hub National}}\\[6pt]
    Pays $\bullet$ Cycles $\bullet$ Classes\\
    Cours $\bullet$ Sections $\bullet$ Domaines\\
    Établissements $\bullet$ Années\\
    Trimestres $\bullet$ Périodes $\bullet$ Mentions\\
    Gestion des Utilisateurs
  };
  
  % Spoke
  \node[block, fill=spokecolor!15, draw=spokecolor!60, right=3cm of hub] (spoke) {
    \textbf{\large MonEcole}\\[4pt]
    \textcolor{spokecolor}{\textbf{Spoke — École}}\\[6pt]
    Inscriptions $\bullet$ Évaluations\\
    Notes $\bullet$ Délibérations\\
    Bulletins PDF $\bullet$ Personnel\\
    Recouvrement $\bullet$ Horaires\\
    Bibliothèque
  };
  
  % Flèches bidirectionnelles
  \draw[arrow] ([yshift=8pt]hub.east) -- ([yshift=8pt]spoke.west) 
    node[midway, above, font=\scriptsize\itshape] {Données de référence};
  \draw[arrow] ([yshift=-8pt]spoke.west) -- ([yshift=-8pt]hub.east) 
    node[midway, below, font=\scriptsize\itshape] {Rapports agrégés};
\end{tikzpicture}
\end{center}

\section{Chiffres Clés}

\begin{center}
\begin{tabular}{>{\bfseries}l r}
\toprule
Indicateur & Valeur \\
\midrule
Pays supportés & 11 (Afrique de l'Est) \\
Modèles Django (Hub) & 25+ \\
Modèles Django (Spoke) & 30+ \\
Endpoints API & 80+ \\
Niveaux hiérarchiques & Dynamiques (0 à N+1) \\
Bases de données & 2+ (Hub + Spokes) \\
Formats de bulletin & 3 (Primaire, Secondaire, Supérieur) \\
\bottomrule
\end{tabular}
\end{center}


% ============================================================
% CHAPITRE 2 : CONTEXTE ET PROBLÉMATIQUE
% ============================================================
\chapter{Contexte et Problématique}

\section{L'État de l'Éducation en Afrique de l'Est}

L'Afrique de l'Est et Centrale regroupe plus de 400 millions d'habitants répartis dans 11 pays ayant chacun un système éducatif distinct. La gestion scolaire dans ces pays est confrontée à plusieurs défis structurels majeurs.

\subsection{Fragmentation des Outils}

La majorité des établissements scolaires utilisent des outils disparates :

\begin{itemize}
  \item \textbf{Registres papier} pour les inscriptions et les notes
  \item \textbf{Fichiers Excel} individuels sans standardisation
  \item \textbf{Logiciels propriétaires} incompatibles entre eux
  \item \textbf{Aucune remontée de données} vers les autorités éducatives
\end{itemize}

\subsection{Absence de Standards Numériques}

Chaque pays possède sa propre nomenclature éducative :

\begin{center}
\begin{tabularx}{\textwidth}{l X X}
\toprule
\textbf{Pays} & \textbf{Structure Pédagogique} & \textbf{Cycles} \\
\midrule
RD Congo & DPE → Sous-Division → École & Primaire, Secondaire, Humanités/CFP \\
Burundi & DPE → DCE → Canton → École & Fondamental, Post-Fondamental \\
Rwanda & District → Secteur → École & Primary, O-Level, A-Level \\
\bottomrule
\end{tabularx}
\end{center}

\subsection{Problèmes Concrets}

\begin{warnbox}[Problèmes Identifiés]
\begin{enumerate}
  \item \textbf{Erreurs de calcul} : Les moyennes et pourcentages sont calculés manuellement, engendrant des erreurs systématiques dans les bulletins.
  \item \textbf{Doublons de données} : Un même élève peut être inscrit dans plusieurs systèmes sans détection.
  \item \textbf{Perte de données} : Les registres papier sont vulnérables aux incendies, inondations et vols.
  \item \textbf{Impossibilité de reporting} : Aucun outil ne permet d'agréger les données au niveau provincial ou national.
  \item \textbf{Coûts cachés} : Le temps passé à reproduire manuellement les bulletins (10-20 min/élève) représente des centaines d'heures par an.
\end{enumerate}
\end{warnbox}

\section{La Réponse : Une Architecture Intégrée}

Face à ces défis, l'écosystème eSchool propose une réponse architecturale en deux couches :

\begin{center}
\begin{tikzpicture}[
  layer/.style={rectangle, rounded corners=6pt, draw, minimum width=14cm, minimum height=2cm, align=center, font=\small},
  arr/.style={-{Stealth[length=6pt]}, thick, gray!60}
]
  \node[layer, fill=hubcolor!12, draw=hubcolor!50] (l1) at (0,3) {
    \textbf{Couche 1 — Gouvernance Nationale} (\texttt{countryStructure})\\
    Standards éducatifs $\bullet$ Hiérarchies pédagogiques et administratives $\bullet$ Catalogues de cours
  };
  
  \node[layer, fill=spokecolor!12, draw=spokecolor!50] (l2) at (0,0) {
    \textbf{Couche 2 — Opérations Scolaires} (\texttt{db\_monecole})\\
    Inscriptions $\bullet$ Notes $\bullet$ Bulletins $\bullet$ Personnel $\bullet$ Finances $\bullet$ Bibliothèque
  };
  
  \draw[arr] (l1.south) -- (l2.north) node[midway, right, font=\scriptsize] {VIEWs MySQL + managed=False};
\end{tikzpicture}
\end{center}


% ============================================================
% CHAPITRE 3 : PRÉSENTATION DE L'ÉCOSYSTÈME
% ============================================================
\chapter{Présentation de l'Écosystème}

\section{Vue d'Ensemble Architecturale}

L'écosystème eSchool repose sur une architecture \textbf{Hub-and-Spoke} (moyeu et rayons), un modèle éprouvé dans les systèmes distribués à grande échelle.

\begin{center}
\begin{tikzpicture}[
  hub/.style={circle, draw=hubcolor, fill=hubcolor!20, minimum size=3cm, font=\small\bfseries, align=center, drop shadow},
  spoke/.style={rectangle, rounded corners=6pt, draw=spokecolor!70, fill=spokecolor!10, minimum width=2.8cm, minimum height=1.8cm, font=\scriptsize, align=center, drop shadow},
  conn/.style={-{Stealth[length=5pt]}, thick, arrowcolor!70}
]
  % Hub central
  \node[hub] (H) {\textcolor{hubcolor}{eSchool}\\Structure\\(\texttt{Hub})};
  
  % Spokes
  \node[spoke] (S1) at (45:5cm) {École Alpha\\Bujumbura\\\texttt{db\_alpha}};
  \node[spoke] (S2) at (105:5cm) {École Beta\\Kinshasa\\\texttt{db\_beta}};
  \node[spoke] (S3) at (165:5cm) {Lycée Gamma\\Kigali\\\texttt{db\_gamma}};
  \node[spoke] (S4) at (225:5cm) {Collège Delta\\Nairobi\\\texttt{db\_delta}};
  \node[spoke] (S5) at (285:5cm) {Institut Epsilon\\Kampala\\\texttt{db\_epsilon}};
  \node[spoke] (S6) at (345:5cm) {École Zeta\\Gitega\\\texttt{db\_zeta}};
  
  % Connexions
  \foreach \s in {S1,S2,S3,S4,S5,S6} {
    \draw[conn] (H) -- (\s);
  }
  
  % Légende
  \node[font=\scriptsize\itshape, text=gray] at (0,-4) {Chaque spoke partage les données de référence du Hub via des VIEWs MySQL cross-BDD};
\end{tikzpicture}
\end{center}

\section{eSchool Structure — Le Hub National}

\texttt{eSchoolStructure} est l'application Django qui gère l'ensemble des données de référence d'un pays. Déployée sur \texttt{https://eschool-rdc.pro/}, elle est le point d'entrée pour les administrateurs nationaux.

\subsection{Responsabilités du Hub}

\begin{enumerate}
  \item \textbf{Structuration du Pays} — Définition des niveaux pédagogiques (DPE → Sous-Division → École) et administratifs (Province → Commune → Colline)
  \item \textbf{Catalogue Éducatif} — Cycles, sections, classes, domaines de cours, catalogue de matières
  \item \textbf{Gestion des Établissements} — Enregistrement, rattachement hiérarchique, configuration annuelle
  \item \textbf{Données Temporelles} — Années scolaires, trimestres, périodes, états (ouvert/clôturé)
  \item \textbf{Tables de Référence} — Sessions d'évaluation, mentions, barèmes
  \item \textbf{Administration des Utilisateurs} — Hiérarchie dynamique, OTP, habilitations
\end{enumerate}

\subsection{Base de Données : \texttt{countryStructure}}

\begin{center}
\begin{tikzpicture}[
  table/.style={rectangle, rounded corners=4pt, draw=hubcolor!60, fill=hubcolor!8, minimum width=3.2cm, minimum height=0.8cm, font=\scriptsize, align=center},
  fk/.style={-{Stealth[length=4pt]}, thin, gray!50}
]
  % Niveau 1 : Pays
  \node[table, fill=hubcolor!20] (pays) at (0,0) {\textbf{pays}};
  
  % Niveau 2 : Structures
  \node[table] (pedtype) at (-5,-1.5) {pedagogicStructuresTypes};
  \node[table] (admtype) at (-1.5,-1.5) {administrativeStructuresTypes};
  \node[table] (cycle) at (1.8,-1.5) {cycles};
  \node[table] (regime) at (5,-1.5) {regimes};
  
  % Niveau 3
  \node[table] (pedinst) at (-5,-3) {pedagogicStructures};
  \node[table] (adminst) at (-1.5,-3) {administrativeStructures};
  \node[table] (section) at (1.8,-3) {sections};
  \node[table] (classe) at (4.5,-3) {classes};
  
  % Niveau 4
  \node[table] (domaine) at (-3,-4.5) {domaines};
  \node[table] (cours) at (0,-4.5) {cours};
  \node[table] (coursannee) at (3,-4.5) {cours\_annee};
  \node[table] (etab) at (6,-4.5) {etablissements};
  
  % Niveau 5
  \node[table] (annee) at (-4,-6) {annees};
  \node[table] (etabannee) at (0,-6) {etab\_annees};
  \node[table] (etabclasse) at (4,-6) {etab\_annees\_classes};
  
  % Niveau 6
  \node[table] (trimestre) at (-2,-7.5) {etab\_annees\_trimestres};
  \node[table] (periode) at (2.5,-7.5) {etab\_annees\_periodes};
  
  % FK arrows
  \foreach \child in {pedtype,admtype,cycle,regime} { \draw[fk] (\child) -- (pays); }
  \draw[fk] (classe) -- (cycle);
  \draw[fk] (cours) -- (classe);
  \draw[fk] (cours) -- (domaine);
  \draw[fk] (coursannee) -- (cours);
  \draw[fk] (coursannee) -- (annee);
  \draw[fk] (etab) -- (pays);
  \draw[fk] (etabannee) -- (etab);
  \draw[fk] (etabannee) -- (annee);
  \draw[fk] (etabclasse) -- (etabannee);
  \draw[fk] (trimestre) -- (etabannee);
  \draw[fk] (periode) -- (trimestre);
\end{tikzpicture}
\end{center}

\section{MonEcole — Le Spoke Scolaire}

\texttt{MonEcole} est l'application Django déployée pour chaque établissement scolaire. Elle consomme les données de référence du Hub et gère les opérations quotidiennes.

\subsection{Responsabilités du Spoke}

\begin{enumerate}
  \item \textbf{Gestion des Élèves} — Inscription, fiche élève complète, suivi de cohorte
  \item \textbf{Évaluations et Notes} — Saisie par période, types de notes (interrogation, examen, TP), maxima configurables
  \item \textbf{Délibérations} — Périodiques, trimestrielles, annuelles, examens, repêchage
  \item \textbf{Bulletins PDF} — Génération automatisée via ReportLab, templates par cycle
  \item \textbf{Personnel} — Enseignants, administratifs, attribution des cours
  \item \textbf{Recouvrement} — Rubriques de frais, paiements, dérogations
  \item \textbf{Bibliothèque} — Armoires, compartiments, livres, emprunts
  \item \textbf{Horaires} — Planification hebdomadaire par classe et enseignant
\end{enumerate}

\subsection{Modèle de Données Simplifié}

\begin{center}
\begin{tikzpicture}[
  tbl/.style={rectangle, rounded corners=4pt, draw=spokecolor!60, fill=spokecolor!8, minimum width=2.6cm, minimum height=0.7cm, font=\scriptsize, align=center},
  viewtbl/.style={rectangle, rounded corners=4pt, draw=orange!60, fill=orange!8, minimum width=2.6cm, minimum height=0.7cm, font=\scriptsize, align=center, dashed},
  fk/.style={-{Stealth[length=4pt]}, thin, gray!50}
]
  % VIEWs (from Hub)
  \node[viewtbl] (vclasse) at (-3,0) {classes (VIEW)};
  \node[viewtbl] (vcours) at (0,0) {cours (VIEW)};
  \node[viewtbl] (vtrim) at (3,0) {trimestres (VIEW)};
  
  % Tables locales
  \node[tbl] (campus) at (-4.5,-1.5) {campus};
  \node[tbl] (annee) at (-1.5,-1.5) {annee};
  \node[tbl] (cycleact) at (1.5,-1.5) {cycle\_actif};
  \node[tbl] (classeact) at (4.5,-1.5) {classe\_active};
  
  \node[tbl] (eleve) at (-4,-3) {eleve};
  \node[tbl] (inscr) at (-1,-3) {inscription};
  \node[tbl] (eval) at (2,-3) {evaluation};
  \node[tbl] (note) at (5,-3) {eleve\_note};
  
  \node[tbl] (delib) at (-2,-4.5) {délibérations};
  \node[tbl] (pers) at (1.5,-4.5) {personnel};
  \node[tbl] (paie) at (4.5,-4.5) {paiements};
  
  % Liens
  \draw[fk] (inscr) -- (eleve);
  \draw[fk] (inscr) -- (classeact);
  \draw[fk] (note) -- (eval);
  \draw[fk] (note) -- (eleve);
  \draw[fk] (classeact) -- (cycleact);
  \draw[fk] (cycleact) -- (campus);
  
  % Légende
  \node[font=\scriptsize\itshape, text=orange!70!black] at (0,0.8) {Tables en pointillés = VIEWs MySQL pointant vers le Hub};
\end{tikzpicture}
\end{center}


% ============================================================
% CHAPITRE 4 : ARCHITECTURE FONCTIONNELLE
% ============================================================
\chapter{Architecture Fonctionnelle}

\section{Le Modèle Hub-and-Spoke en Détail}

L'architecture Hub-and-Spoke est le fondement technique de l'écosystème. Elle résout un problème fondamental : \textbf{comment standardiser les données éducatives à l'échelle nationale tout en préservant l'autonomie opérationnelle de chaque école ?}

\begin{infobox}[Principe Fondamental]
Le Hub \textbf{définit} (structures, cours, calendrier). Le Spoke \textbf{exécute} (inscriptions, notes, bulletins). La synchronisation est \textbf{transparente} via des VIEWs MySQL cross-bases de données.
\end{infobox}

\subsection{Flux de Données}

\begin{center}
\begin{tikzpicture}[
  process/.style={rectangle, rounded corners=6pt, draw, minimum width=4cm, minimum height=1.5cm, align=center, font=\small, drop shadow={shadow xshift=1pt, shadow yshift=-1pt}},
  data/.style={cylinder, draw, shape border rotate=90, minimum width=1.5cm, minimum height=1cm, font=\scriptsize, align=center, aspect=0.3},
  arr/.style={-{Stealth[length=6pt]}, thick},
  darr/.style={-{Stealth[length=6pt]}, thick, dashed}
]
  % Acteurs
  \node[process, fill=eschoolpurple!15, draw=eschoolpurple!60] (admin) at (0,4) {
    \textbf{Admin National}\\
    Configure le pays
  };
  
  \node[process, fill=hubcolor!15, draw=hubcolor!60] (hub) at (0,1.5) {
    \textbf{eSchool Structure}\\
    \texttt{countryStructure}
  };
  
  \node[data, fill=hubcolor!10, draw=hubcolor!50] (dbhub) at (5,1.5) {\texttt{country}\\\texttt{Structure}};
  
  \node[process, fill=spokecolor!15, draw=spokecolor!60] (spoke1) at (-4,-1.5) {
    \textbf{École A}\\
    \texttt{db\_ecoleA}
  };
  \node[process, fill=spokecolor!15, draw=spokecolor!60] (spoke2) at (0,-1.5) {
    \textbf{École B}\\
    \texttt{db\_ecoleB}
  };
  \node[process, fill=spokecolor!15, draw=spokecolor!60] (spoke3) at (4,-1.5) {
    \textbf{École C}\\
    \texttt{db\_ecoleC}
  };
  
  \node[process, fill=eschoolamber!15, draw=eschoolamber!60] (dir) at (0,-4) {
    \textbf{Directeur / Enseignant}\\
    Gère son école
  };
  
  % Flèches
  \draw[arr, eschoolpurple!70] (admin) -- (hub);
  \draw[arr, hubcolor!70] (hub) -- (dbhub);
  \draw[arr, arrowcolor] (dbhub.south) -- ++(0,-0.5) -| (spoke1.north);
  \draw[arr, arrowcolor] (dbhub.south) -- ++(0,-0.5) -| (spoke2.north);
  \draw[arr, arrowcolor] (dbhub.south) -- ++(0,-0.5) -| (spoke3.north);
  \draw[arr, eschoolamber!70] (dir) -- (spoke2);
  
  % Labels
  \node[font=\scriptsize\itshape, text=arrowcolor] at (5,-0.2) {VIEWs MySQL};
\end{tikzpicture}
\end{center}

\subsection{Mécanismes de Synchronisation}

La synchronisation entre le Hub et les Spokes repose sur trois mécanismes complémentaires :

\begin{enumerate}
  \item \textbf{VIEWs MySQL Cross-BDD} — Des vues SQL dans la base spoke pointent directement vers les tables du Hub. L'ORM Django ne voit aucune différence.

  \item \textbf{Modèles \texttt{managed=False}} — Les modèles Django correspondant aux tables du Hub sont déclarés avec \texttt{managed = False} dans le Spoke, empêchant Django de modifier leur schéma.

  \item \textbf{DatabaseRouter} — Un routeur Django dirige les requêtes ORM vers la bonne base de données selon le modèle interrogé.
\end{enumerate}

\begin{center}
\begin{tikzpicture}[
  box/.style={rectangle, rounded corners=4pt, draw, minimum width=4.5cm, minimum height=1.2cm, align=center, font=\small},
  arr/.style={-{Stealth[length=5pt]}, thick}
]
  \node[box, fill=spokecolor!10, draw=spokecolor!50] (django) at (0,0) {\textbf{Django ORM}\\MonEcole};
  \node[box, fill=orange!10, draw=orange!50] (router) at (5,0) {\textbf{DatabaseRouter}\\Aiguillage};
  \node[box, fill=spokecolor!15, draw=spokecolor!60] (dblocal) at (10,1) {\texttt{db\_monecole}\\Tables locales};
  \node[box, fill=hubcolor!15, draw=hubcolor!60] (dbhub) at (10,-1) {\texttt{countryStructure}\\Tables Hub};
  
  \draw[arr] (django) -- (router) node[midway, above, font=\scriptsize] {Requête};
  \draw[arr, spokecolor] (router) -- (dblocal) node[midway, above, font=\scriptsize] {Élèves, Notes...};
  \draw[arr, hubcolor] (router) -- (dbhub) node[midway, below, font=\scriptsize] {Cours, Trimestres...};
\end{tikzpicture}
\end{center}

\section{Identification du Tenant}

Chaque école (tenant) est identifiée par un \texttt{id\_etablissement} stocké dans la table \texttt{campus} du Spoke. Cette clé logique permet de résoudre la chaîne suivante :

\begin{center}
\begin{tikzpicture}[
  step/.style={rectangle, rounded corners=4pt, draw=arrowcolor!50, fill=arrowcolor!8, minimum width=2.8cm, minimum height=1cm, align=center, font=\scriptsize},
  arr/.style={-{Stealth[length=5pt]}, thick, arrowcolor!60}
]
  \node[step] (s1) at (0,0) {Campus local\\(Spoke)};
  \node[step] (s2) at (3.5,0) {Établissement\\(Hub)};
  \node[step] (s3) at (7,0) {Année activée\\(Hub)};
  \node[step] (s4) at (10.5,0) {Trimestres\\Périodes};
  
  \draw[arr] (s1) -- (s2) node[midway, above, font=\tiny] {\texttt{id\_etab}};
  \draw[arr] (s2) -- (s3) node[midway, above, font=\tiny] {\texttt{etab\_annees}};
  \draw[arr] (s3) -- (s4) node[midway, above, font=\tiny] {\texttt{etab\_trim}};
\end{tikzpicture}
\end{center}

\section{Résolution Multi-Pays par Sous-Domaine}

Le système identifie automatiquement le pays via le sous-domaine de l'URL :

\begin{center}
\begin{tabular}{l l l}
\toprule
\textbf{URL} & \textbf{Sous-domaine} & \textbf{Pays résolu} \\
\midrule
\texttt{rdc.monecole.pro} & \texttt{rdc} & RD Congo \\
\texttt{bdi.monecole.pro} & \texttt{bdi} & Burundi \\
\texttt{rw.monecole.pro} & \texttt{rw} & Rwanda \\
\texttt{eschool-rdc.pro} & (domaine principal) & RD Congo \\
\bottomrule
\end{tabular}
\end{center}


% ============================================================
% CHAPITRE 5 : MODULES FONCTIONNELS
% ============================================================
\chapter{Modules Fonctionnels}

L'écosystème eSchool comprend \textbf{12 modules fonctionnels} répartis entre le Hub et les Spokes.

\begin{center}
\begin{tikzpicture}[
  mod/.style={rectangle, rounded corners=6pt, draw, minimum width=3cm, minimum height=1cm, align=center, font=\scriptsize\bfseries, drop shadow={shadow xshift=0.5pt, shadow yshift=-0.5pt}},
]
  % Hub modules (bleu)
  \node[mod, fill=hubcolor!15, draw=hubcolor!50] (m1) at (-5,4) {Structuration\\Pays};
  \node[mod, fill=hubcolor!15, draw=hubcolor!50] (m2) at (-1.5,4) {Cycles \&\\Sections};
  \node[mod, fill=hubcolor!15, draw=hubcolor!50] (m3) at (2,4) {Programmes\\(Cours)};
  \node[mod, fill=hubcolor!15, draw=hubcolor!50] (m4) at (5.5,4) {Établissements};
  
  \node[mod, fill=hubcolor!15, draw=hubcolor!50] (m5) at (-3.5,2.5) {Config.\\Annuelle};
  \node[mod, fill=hubcolor!15, draw=hubcolor!50] (m6) at (0,2.5) {Utilisateurs\\Admin};
  \node[mod, fill=hubcolor!15, draw=hubcolor!50] (m7) at (3.5,2.5) {Références\\(Mentions...)};
  
  % Spoke modules (vert)
  \node[mod, fill=spokecolor!15, draw=spokecolor!50] (m8) at (-5,0.5) {Inscriptions};
  \node[mod, fill=spokecolor!15, draw=spokecolor!50] (m9) at (-1.5,0.5) {Évaluations\\Notes};
  \node[mod, fill=spokecolor!15, draw=spokecolor!50] (m10) at (2,0.5) {Délibérations};
  \node[mod, fill=spokecolor!15, draw=spokecolor!50] (m11) at (5.5,0.5) {Bulletins\\PDF};
  
  \node[mod, fill=spokecolor!15, draw=spokecolor!50] (m12) at (-3.5,-1) {Personnel};
  \node[mod, fill=spokecolor!15, draw=spokecolor!50] (m13) at (0,-1) {Recouvrement};
  \node[mod, fill=spokecolor!15, draw=spokecolor!50] (m14) at (3.5,-1) {Bibliothèque};
  
  % Labels
  \node[font=\small\bfseries, text=hubcolor] at (-6.5,3.4) {\rotatebox{90}{HUB}};
  \node[font=\small\bfseries, text=spokecolor] at (-6.5,-0.2) {\rotatebox{90}{SPOKE}};
  
  % Séparation
  \draw[dashed, gray!40, thick] (-7,1.6) -- (7,1.6);
\end{tikzpicture}
\end{center}

\section{Module 1 : Structuration du Pays}

Ce module permet de définir la hiérarchie pédagogique et administrative d'un pays. Le nombre de niveaux est dynamique — défini par \texttt{nLevelsStructuraux} et \texttt{nLevelsAdministratifs}.

\begin{successbox}[Exemple : RD Congo]
\textbf{Structure Pédagogique (5 niveaux)} : Division Provinciale de l'Éducation → Sous-Division → Inspection → École → Classe

\textbf{Structure Administrative (4 niveaux)} : Province → Territoire → Secteur → Village
\end{successbox}

\section{Module 2 : Cycles, Sections et Classes}

Les cycles d'enseignement (Primaire, Secondaire, Humanités...) sont définis par pays avec :
\begin{itemize}
  \item Un \textbf{ordre} hiérarchique
  \item Une \textbf{durée} en années
  \item Un flag \texttt{hasSections} pour les cycles nécessitant des sections (Scientifique, Littéraire...)
  \item Des \textbf{classes} auto-générées selon la durée du cycle
\end{itemize}

\section{Module 3 : Programmes et Cours}

Le catalogue des cours suit un modèle en deux niveaux :

\begin{center}
\begin{tikzpicture}[
  box/.style={rectangle, rounded corners=4pt, draw, minimum width=4cm, minimum height=1.5cm, align=center, font=\small},
  arr/.style={-{Stealth[length=5pt]}, thick}
]
  \node[box, fill=hubcolor!10, draw=hubcolor!50] (cat) at (0,0) {\textbf{Catalogue (Cours)}\\Identité permanente\\Code + Nom + Domaine};
  \node[box, fill=eschoolamber!10, draw=eschoolamber!50] (ann) at (6,0) {\textbf{Config Annuelle}\\(\texttt{CoursAnnee})\\Pondération, heures, crédits};
  
  \draw[arr] (cat) -- (ann) node[midway, above, font=\scriptsize] {1:N par année};
\end{tikzpicture}
\end{center}

\section{Module 4 : Évaluations et Notes}

Le système d'évaluation supporte :
\begin{itemize}
  \item \textbf{5 types de notes} : Interrogation, Travail, Examen, TP, TPE
  \item \textbf{Sessions} : Ordinaire et Repêchage
  \item \textbf{Périodes} : P1..P6 regroupées en Trimestres T1..T3
  \item \textbf{Maxima configurables} par cours et par année
\end{itemize}

\section{Module 5 : Délibérations}

Cinq types de délibérations sont supportés :
\begin{enumerate}
  \item Délibération \textbf{Périodique} — Par période (P1, P2...)
  \item Délibération \textbf{Trimestrielle} — Par trimestre
  \item Délibération \textbf{Annuelle} — Fin d'année avec finalités (Réussi, Échoué, Ajourné)
  \item Délibération \textbf{d'Examen} — Sessions de fin de cycle
  \item Délibération \textbf{de Repêchage} — Seconde chance
\end{enumerate}

\section{Module 6 : Bulletins PDF}

Les bulletins sont générés automatiquement via \textbf{ReportLab} avec des templates adaptés à chaque cycle :

\begin{center}
\begin{tabular}{l l l}
\toprule
\textbf{Cycle} & \textbf{Template} & \textbf{Particularités} \\
\midrule
Primaire & \texttt{bulletin\_primaire.py} & Format simplifié, sans crédits \\
Secondaire & \texttt{bulletin\_secondaire.py} & Grille complète avec domaines \\
Supérieur & \texttt{bulletin\_superieur.py} & Crédits, CM/TD/TP, ECTS \\
\bottomrule
\end{tabular}
\end{center}

\section{Module 7 : Bibliothèque}

Le module bibliothèque gère l'inventaire physique et les emprunts :
\begin{itemize}
  \item \textbf{Armoires} → \textbf{Compartiments} → \textbf{Livres} → \textbf{Exemplaires}
  \item Système d'emprunt avec profils et limites
  \item Statistiques de fréquentation
\end{itemize}


% ============================================================
% CHAPITRE 6 : MODÈLE MULTI-PAYS
% ============================================================
\chapter{Modèle Multi-Pays}

\section{Pays Supportés}

Le système est conçu pour supporter les 11 pays d'Afrique de l'Est et Centrale :

\begin{center}
\begin{tabular}{c l l l}
\toprule
\textbf{Sigle} & \textbf{Pays} & \textbf{Domaine} & \textbf{Statut} \\
\midrule
CD & RD Congo & eschool-rdc.pro & \textcolor{green!60!black}{\textbf{Actif}} \\
BDI & Burundi & bdi.monecole.pro & \textcolor{orange}{\textbf{Pilote}} \\
RW & Rwanda & rw.monecole.pro & Planifié \\
KE & Kenya & ke.monecole.pro & Planifié \\
TZ & Tanzanie & tz.monecole.pro & Planifié \\
UG & Ouganda & ug.monecole.pro & Planifié \\
SS & Soudan du Sud & ss.monecole.pro & Planifié \\
ET & Éthiopie & et.monecole.pro & Planifié \\
SO & Somalie & so.monecole.pro & Planifié \\
DJ & Djibouti & dj.monecole.pro & Planifié \\
ER & Érythrée & er.monecole.pro & Planifié \\
\bottomrule
\end{tabular}
\end{center}

\section{Adaptabilité par Pays}

Chaque pays peut configurer indépendamment :
\begin{itemize}
  \item Le \textbf{nombre de niveaux} pédagogiques et administratifs
  \item Les \textbf{noms des niveaux} (DPE, Sous-Division... ou District, Secteur...)
  \item Les \textbf{cycles d'enseignement} et leur durée
  \item Les \textbf{sections} (Scientifique, Littéraire, Technique...)
  \item Les \textbf{régimes} (Public, Privé, Conventionné)
  \item Les \textbf{programmes} d'études (National, Belge, Français...)
  \item Les \textbf{barèmes de notation} (Mentions)
\end{itemize}


% ============================================================
% CHAPITRE 7 : SÉCURITÉ
% ============================================================
\chapter{Sécurité et Gestion des Accès}

\section{Hiérarchie Dynamique}

Le modèle \texttt{AdminUser} implémente une hiérarchie à N+2 niveaux déterminée dynamiquement :

\begin{center}
\begin{tikzpicture}[
  level/.style={rectangle, rounded corners=4pt, draw, minimum width=5cm, minimum height=1cm, align=center, font=\small},
  arr/.style={-{Stealth[length=5pt]}, thick, gray!60}
]
  \node[level, fill=red!10, draw=red!40] (l0) at (0,0) {\textbf{Niveau 0} — National\\Accès total au pays};
  \node[level, fill=orange!10, draw=orange!40] (l1) at (0,-1.8) {\textbf{Niveau 1..N} — Structure Pédagogique\\Accès à une branche};
  \node[level, fill=spokecolor!10, draw=spokecolor!40] (l2) at (0,-3.6) {\textbf{Niveau N+1} — Établissement\\Accès à une école};
  
  \draw[arr] (l0) -- (l1) node[midway, right, font=\scriptsize] {Crée des admins de niveau 1};
  \draw[arr] (l1) -- (l2) node[midway, right, font=\scriptsize] {Crée des admins d'école};
\end{tikzpicture}
\end{center}

\section{Authentification Multi-Étapes}

\begin{enumerate}
  \item \textbf{Découverte d'identité} — L'utilisateur saisit son email, le système vérifie son existence
  \item \textbf{Branchement} :
  \begin{itemize}
    \item Utilisateur vérifié avec mot de passe → Formulaire de connexion
    \item Utilisateur non vérifié → Challenge OTP (Email ou SMS)
  \end{itemize}
  \item \textbf{Vérification OTP} — Code à 6 chiffres avec expiration de 10 minutes, 3 tentatives max
  \item \textbf{Création de mot de passe} — Pour les nouveaux comptes après vérification OTP
  \item \textbf{Session sécurisée} — Timeout d'inactivité de 10 minutes, PBKDF2-SHA256
\end{enumerate}


% ============================================================
% CHAPITRE 8 : FEUILLE DE ROUTE
% ============================================================
\chapter{Feuille de Route}

\section{Phases de Développement}

\begin{center}
\begin{tikzpicture}[
  phase/.style={rectangle, rounded corners=6pt, draw, minimum width=3.5cm, minimum height=2.5cm, align=center, font=\small, drop shadow},
  arr/.style={-{Stealth[length=6pt]}, ultra thick}
]
  \node[phase, fill=spokecolor!15, draw=spokecolor!50] (p1) at (0,0) {
    \textbf{Phase 1}\\[2pt]
    \textcolor{spokecolor}{\textbf{MVP}}\\[4pt]
    {\scriptsize Hub + 1 Spoke}\\
    {\scriptsize RDC opérationnel}
  };
  
  \node[phase, fill=hubcolor!15, draw=hubcolor!50] (p2) at (4.5,0) {
    \textbf{Phase 2}\\[2pt]
    \textcolor{hubcolor}{\textbf{Multi-Tenant}}\\[4pt]
    {\scriptsize Provisioning auto}\\
    {\scriptsize Multi-pays}
  };
  
  \node[phase, fill=eschoolpurple!15, draw=eschoolpurple!50] (p3) at (9,0) {
    \textbf{Phase 3}\\[2pt]
    \textcolor{eschoolpurple}{\textbf{Mobile}}\\[4pt]
    {\scriptsize App parents/élèves}\\
    {\scriptsize Notifications push}
  };
  
  \node[phase, fill=eschoolamber!15, draw=eschoolamber!50] (p4) at (13.5,0) {
    \textbf{Phase 4}\\[2pt]
    \textcolor{eschoolamber}{\textbf{Analytics}}\\[4pt]
    {\scriptsize Tableaux de bord}\\
    {\scriptsize IA prédictive}
  };
  
  \draw[arr, spokecolor!50] (p1) -- (p2);
  \draw[arr, hubcolor!50] (p2) -- (p3);
  \draw[arr, eschoolpurple!50] (p3) -- (p4);
\end{tikzpicture}
\end{center}

% ============================================================
% CHAPITRE 9 : IMPACT
% ============================================================
\chapter{Impact Attendu}

\section{Métriques d'Impact}

\begin{center}
\begin{tabularx}{\textwidth}{X c c}
\toprule
\textbf{Indicateur} & \textbf{Avant eSchool} & \textbf{Avec eSchool} \\
\midrule
Temps de génération d'un bulletin & 15--20 min & < 1 seconde \\
Taux d'erreurs dans les moyennes & 5--10\% & 0\% \\
Remontée de données au ministère & Aucune & Temps réel \\
Temps d'inscription d'un élève & 20 min & 3 min \\
Coût annuel par école (estimé) & Variable & \$50--100/an \\
Risque de perte de données & Élevé & Quasi-nul \\
\bottomrule
\end{tabularx}
\end{center}

% ============================================================
% ANNEXES
% ============================================================
\appendix
\chapter{Glossaire}

\begin{description}
  \item[Hub] Base de données centralisée \texttt{countryStructure} contenant les données de référence nationales
  \item[Spoke] Base de données locale \texttt{db\_monecole} propre à chaque établissement scolaire
  \item[VIEW] Vue MySQL permettant de lire des données d'une autre base comme si elles étaient locales
  \item[ORM] Object-Relational Mapping — Couche d'abstraction Django pour les requêtes SQL
  \item[OTP] One-Time Password — Code à usage unique pour la vérification d'identité
  \item[Tenant] Entité isolée dans un système multi-locataires (ici : un établissement)
  \item[Cycle] Niveau d'enseignement (Primaire, Secondaire, Humanités...)
  \item[Section] Filière au sein d'un cycle (Scientifique, Littéraire, Technique...)
  \item[Domaine] Regroupement thématique de cours (Sciences, Langues, Arts...)
  \item[CoursAnnee] Configuration annuelle d'un cours (pondération, heures, crédits)
  \item[Délibération] Processus de validation des résultats scolaires
  \item[Mention] Appréciation qualitative basée sur un pourcentage (Très Bien, Bien, Assez Bien...)
\end{description}

\end{document}
